\ifdefined\iscombined
	\lecturetitle{Actions and Labels}
	\setcounter{section}{6}
\else
\documentclass{article}
\usepackage{etoolbox}
\usepackage{enumitem}
\usepackage{amsthm}
\usepackage{comment}
\usepackage{amsmath}
\usepackage{amsfonts}
\usepackage{sectsty}
\usepackage{hyperref}
\usepackage{fancyhdr}
\usepackage[top=1in,bottom=1in,left=1in,right=1in]{geometry}
\usepackage{float}
\usepackage{graphicx}
\usepackage{tabularx}
\usepackage{mathtools}

\title{Lecture 6 - Actions and Labels}
\author{Matthew Farah, \href{mailto:farahm11@mcmaster.ca}{$\langle$farahm11@mcmaster.ca$\rangle$}, 400468251}
\date{\today}

\hypersetup{
	colorlinks=true,
	linkcolor=blue,
	filecolor=magenta,
	urlcolor=blue,
	citecolor=blue,
	anchorcolor=blue
}

\fancyhead{}
\fancyhead[L]{Matthew Farah, \href{mailto:farahm11@mcmaster.ca}{$\langle$farahm11@mcmaster.ca$\rangle$}, 400468251}
\fancyhead[R]{Lecture 6 - Actions and Labels}

\newcounter{example}
\newcounter{partctr}[example]

\newenvironment{example}{
	\refstepcounter{example}
	\setcounter{partctr}{0}
	\vspace{1em}
	\large\noindent\textbf{Ex.}\vspace{0.5em}
}{\vspace{1.5em}}

\newenvironment{parts}{
	\begin{enumerate}[
		label=\textbf{(\alph*)},
		leftmargin=3.3em,
		itemsep=0.4em,
		before=\large
	]
	\setcounter{enumi}{\value{partctr}}
}{
	\setcounter{partctr}{\value{enumi}}
	\end{enumerate}
}

\newenvironment{solution}{
	\vspace{0.3cm}\par\noindent\large\textbf{Solution:}\par\vspace{0.3em}
}{\vspace{0.8em}}

\renewcommand{\thesubsection}{\thesection.\alph{subsection}}
\renewcommand{\thesubsubsection}{\thesection.\alph{subsection}.\Roman{subsubsection}}
\makeatletter

\sectionfont{\fontsize{12}{12}\bfseries}
\subsectionfont{\fontsize{10}{12}\normalfont}
\subsubsectionfont{\fontsize{10}{12}\normalfont}

\newcommand{\parts@seriesname}{parts@\theexample}

\begin{document}

\maketitle
\pagestyle{fancy}

\fi

\begin{center}
\textbf{These notes are a bit of a stub because I don't totally understand what this lecture was talking about, I'll try to update with better notes in future.}\\
\end{center}

\begin{itemize}
	\item \emph{Action relabelling} is the act of renaming actions in a process.
	\begin{itemize}
		\item This is done to force concurrently running processes to execute certain actions simultaneously, even if they have different names.
	\end{itemize}
	\item Actions can be relabelled with a prefix.
	\item \emph{Action hiding} is the act of removing certain actions from the alphabet of a process.
	\begin{itemize}
		\item This is done to prevent actions with the same name from being shared between two concurrently running processes.
		%TODO: WHAT?
		\item Silent actions are relabelled as $\tau$.
	\end{itemize}
	%TODO: WHAT?
	\item \emph{Structural diagrams} are used to represent:
	\begin{enumerate}
		\item Concurrent systems.
		\item Action relabelling.
		\item Action hiding.
	\end{enumerate}
	\item An \emph{active} process is a process which initiates actions.
	\item A \emph{passive} process is a process which responds to actions.
\end{itemize}

\ifdefined\iscombined
	\newpage
\else
	\end{document}
\fi
